\section{Introduction}

The nocturnal Stable Boundary Layer (SBL) remains one of the most persistent challenges in geophysical fluid dynamics due to the highly non-linear, non-equilibrium coupling between turbulent dissipation, surface energy budgets, and mesoscale forcing \citep{nieuwstadt1984, mahrt2014, vanDeWiel2017}. Observational studies consistently report substantial variability in the critical Bulk Richardson number ($Ri_c$) across field campaigns and environmental regimes, despite traditional expectations that it acts as a quasi-universal constant near $0.25$ \citep{Grachev2013, Monahan2015}. While mid-latitude field campaigns such as CASES-99 report collapse thresholds near the classical limit ($Ri \approx 0.2\text{--}0.4$), polar campaigns like SHEBA record active turbulence persisting at stability ratios exceeding unity ($Ri > 1.2$) \citep{poulos2002cases99, Grachev2005}. We refer to this long-standing observational inconsistency as the \emph{Richardson universality problem}.

\subsection{The Paradox of Stability Thresholds}

For decades, the search for a refined stability function has assumed that variability in $Ri_c$ stems from missing physics or measurement noise. Under the classical equilibrium paradigm, a unique critical Richardson number should exist after accounting for measurement uncertainty. However, traditional parameterizations---such as Monin--Obukhov Similarity Theory (MOST) or Mellor--Yamada $K$-theory---rely on an instantaneous local equilibrium assumption, $e = e(S, N^2)$, that erases the system's underlying temporal history and structural hysteresis \citep{MellorYamada1982}. This diagnostic approach forces a smooth, single-valued response that cannot account for the abrupt transitions between weakly stable and very stable regimes. The persistent campaign-to-campaign variability therefore motivates a fundamentally geometric reinterpretation.

\subsection{Thesis: Richardson Thresholds as Projections of Invariant Geometry}

In this paper, we propose that the Richardson threshold problem is naturally explained by projection from a higher-dimensional invariant manifold rather than incompatible physics. We frame the SBL as a singularly perturbed dynamical system governed by Geometric Singular Perturbation Theory (GSPT) \citep{Fenichel1979, Kuehn2015}. Within this framework, our central thesis is that \emph{observed Richardson thresholds are projection-dependent diagnostics of an invariant fold geometry}.

By exploiting the natural separation of timescales---where fast turbulence adjustment ($\tau_e \sim 10^1\text{--}10^2\,\text{s}$) responds to slow surface thermal and geostrophic evolution ($\tau_{\text{slow}} \sim 10^3\text{--}10^4\,\text{s}$)---we demonstrate that the critical threshold is not a static point, but a state-dependent fold locus ($\mathcal{C}_{\mathrm{fold}}$) that deforms dynamically in response to the Surface Energy Balance (SEB) \citep{VanDeWiel2012}. Crucially, observational site constraints ($\Sigma_{\mathrm{site}}$) intersect the invariant critical manifold ($\mathcal{S}_0^+$) transversally ($\mathcal{S}_0^+ \pitchfork \Sigma_{\mathrm{site}}$), yielding a unique campaign-specific trajectory whose coordinate projection generates the observed scalar Richardson threshold:
$$Ri_{\mathrm{obs}} = \Pi_{Ri}\left(\mathcal{C}_{\mathrm{fold}} \cap \Sigma_{\mathrm{site}}\right)$$

\subsection{The ``Fold Illusion'' and Emergent Coupling}

A key conceptual contribution of this work is the distinction between isolated boundary-driven decay and true fold catastrophes. We show that while an isolated atmospheric subsystem exhibits only a monotonic boundary-crossing at the laminar floor ($e = 0$), a genuine S-shaped fold only emerges through the non-linear coupling of the atmosphere to the SEB. We call this phenomenon the ``Fold Illusion,'' and demonstrate that resolving this distinction provides a geometric mechanism for modeling intermittent bursting and nocturnal relaxation oscillations accurately.

\subsection{Objectives and Scope of Part 1}

As the first in a three-part series, this manuscript establishes the mathematical foundations and observational resolution of the GSPT-SBL framework. We define a hierarchy of models---from 2D analytical baselines to 5D thermodynamic systems---and provide a sequence of structural theorems:
\begin{itemize}
    \item \textbf{Theorem 1 (Existence of the Critical Manifold):} Establishing $\mathcal{S}_0^+$ as a smooth embedded submanifold of the fast-slow state space via Fenichel theory.
    \item \textbf{Theorem 2 (Characterization of the Fold Locus):} Defining the geometric boundary $\mathcal{C}_{\mathrm{fold}}$ where normal hyperbolicity breaks down, triggering turbulence collapse.
    \item \textbf{Theorem 3 (Projection of the Fold Geometry):} Formalizing the mapping $\Pi_{Ri}$ from invariant state-space intersections onto campaign-specific scalar diagnostics ($Ri_{\mathrm{fold}}$).
\end{itemize}

Part 1 provides the geometric skeleton required for the data-driven discovery and prognostic parameterizations developed in Parts 2 and 3.
